% \subsection{Summary of the PDR}
% In this section, we will provide a brief summary of the Preliminary Design Review (PDR) for our Rocketry/CanSat mission. The PDR covers a range of topics, including the mission objectives, construction of the CanSat, electrical design, software development, risk mitigation strategies, data analysis plan, and outreach program. We will highlight the key aspects of each of these areas and provide an overview of the overall mission.

% \begin{itemize}[leftmargin=1cm, itemindent=0.25cm, noitemsep, topsep=0pt, label=$\bullet$]
% \item[I.] Introduction
% \begin{itemize}[label=$\bullet$, noitemsep, topsep=2pt]
% \item Details of the CanSat mission to analyze a planet's environment, including air quality, light conditions, wind speeds, terrain type, temperature, pressure, and live video to explore it further.
% \item The team members' roles and responsibilities, including team leader, embedded software developer, mission data analyst, mechanical and structural designer, and hardware developer and designer.
% \item Two primary missions: measuring air pollution and analyzing a potentially habitable planet by taking samples to assess the sustainability of life.
% \end{itemize}
% \item[II.] CanSat description
% \begin{itemize}[label=$\bullet$, noitemsep, topsep=2pt]
% \item The CanSat will measure various data during its descent, including temperature, pressure, space positioning, location, UV radiation, pollution gases, and particulate matter concentration, and will transmit this data in real time to a ground station and store it onboard.
% \item The mechanical structure is made of lightweight materials, including thermoplastic, carbon material, and 3D-printed Polylactic Acid (PLA) with carbon fibers glued with epoxy, a shock-absorbing system, a stabilization wing concept, and a parachute to safely return it to the ground.
% \item The electrical design of a custom PCB for a can that uses multiple sensors and modules, including two GPS modules, a particulate matter sensor, gas sensors, an RF modem, and a sensor fusion module.
% \item The design of the software for the project, which has two main purposes: acquiring and logging data from various sensors.
% \item A custom piece of software is developed to run on the ground station to interpret and visualize telemetry data.
% \item The recovery system for the CanSat consists of a hemispherical parachute, live telemetry (GSM, video, and RF), and a precise localization system based on audio-visual aids.
% \end{itemize}
% \item[III.] Project planning
% \begin{itemize}[label=$\bullet$, noitemsep, topsep=2pt]
% \item The risks associated with the project, include launch vehicle failure, communication failure, power limitations, environmental factors, technical issues, budget constraints, and resource estimation.
% \item The team's steps to mitigate risks, include creating a detailed power budget, conducting endurance tests, developing a prototype CanSat development board, managing a budget, and estimating resources accurately.
% \end{itemize}
% \item[IV.] Data analysis
% \begin{itemize}[label=$\bullet$, noitemsep, topsep=2pt]
% \item The data analysis plan includes various software programs and tools to collect, process, clean, visualize, and analyze the data collected by the CanSat sensors.
% \item The outreach program includes various strategies to promote the CanSat competition and team, including presenting ideas about CanSat to classmates, promoting the competition on social media platforms, organizing online workshops to teach electronics to kids and interested individuals, seeking sponsorships from businesses or organizations, collaborating with local schools to introduce CanSat, and participating in events and fairs related to science and technology.
% \end{itemize}
% \end{itemize}

\subsection{Summary of the PDR}
In this section, we will provide a brief summary of the Preliminary Design Review (PDR) for our CanSat mission. The PDR covers a range of topics, including the mission objectives, construction of the CanSat, electrical design, software development, risk mitigation strategies, data analysis plan, and outreach program. We will highlight the key aspects of each of these areas and provide an overview of the overall mission.

\begin{itemize}[leftmargin=1cm, itemindent=0.25cm, noitemsep, topsep=0pt, label=$\bullet$]
\item[I.] Introduction
\begin{itemize}[label=$\bullet$, noitemsep, topsep=2pt]
\item Overview of the CanSat mission objectives, focusing on environmental monitoring and cosmic ray analysis, including muon detection, UV light intensity measurement, and carbon dioxide/monoxide level assessment.
\item Description of team roles and responsibilities, such as team leader, base operator, CanSat and sensor engineer, data analyst, software developer, and mechanical designer.
\item Emphasis on two primary missions: muon detection to understand cosmic ray interactions, and environmental monitoring to analyze UV radiation and CO levels for health and environmental studies.
\end{itemize}
\item[II.] CanSat description
\begin{itemize}[label=$\bullet$, noitemsep, topsep=2pt]
\item The CanSat's capabilities in measuring environmental data (temperature, pressure, UV radiation, CO$_2$/CO levels), space positioning, and muon counts during its descent, along with real-time data transmission to a ground station and onboard data storage.
\item Mechanical design details, including the use of lightweight and durable materials, a shock-absorbing system, stabilization features, and a recovery parachute.
\item The electrical design encompassing a custom PCB that integrates muon detectors, UV and CO$_2$/CO sensors, GPS modules, and an RF modem for communication.
\item Software development tailored for sensor data acquisition, logging, and ground station telemetry interpretation and visualization.
\item A comprehensive recovery system featuring a parachute, live telemetry, and a precise localization mechanism.
\end{itemize}
\item[III.] Rocket design
\begin{itemize}[label=$\bullet$, noitemsep, topsep=2pt]
\item Design specifications for the rocket's structural components, focusing on the nose cone, fuselage, and fins, and select appropriate materials for construction.
\item Integration and testing the parachute recovery system. Conduct stability simulations to fine-tune the fin design and validate the rocket's flight dynamics.
\end{itemize}
\item[IV.] Project planning
\begin{itemize}[label=$\bullet$, noitemsep, topsep=2pt]
\item Identification of project risks, including payload integrity, sensor accuracy, communication challenges, power constraints, and environmental impacts.
\item Risk mitigation strategies, such as rigorous testing of muon detectors and environmental sensors, power budget optimization, and robust communication protocols.
\end{itemize}
\item[V.] Data analysis
\begin{itemize}[label=$\bullet$, noitemsep, topsep=2pt]
\item The plan for data analysis, using advanced software to process, visualize, and interpret the collected muon, UV, and CO$_2$/CO data.
\item An outreach program aimed at promoting the scientific significance of the mission, engaging with educational institutions, and utilizing social media and public events to disseminate project information and findings.
\end{itemize}
\end{itemize}

